%% Cuerpo de la constitución de Legasoft (LGS-INST-001).
%% Fragmento: no se compila solo. Lo incluyen constitucion-legasoft.tex
%% y el documento de entrega docs/entrega/entrega-fase1.tex.

% =============================================================================
\section{Introducción}

La presente propuesta establece la identidad, el propósito y la organización
inicial de \textbf{Legasoft}, una consultora tecnológica simulada en el marco de
la formación universitaria en Ingeniería de Sistemas. Su creación responde a la
necesidad de articular conocimientos técnicos, gestión organizacional y prácticas
profesionales en una estructura capaz de analizar problemas reales y proponer
soluciones de software viables, documentadas y orientadas a resultados.

La ingeniería de software profesional trasciende la programación de una
aplicación. Comprende un proceso gestionado que inicia con la comunicación con el
cliente y el conocimiento de su contexto, continúa con el levantamiento y la
gestión de requisitos, la planificación, el diseño arquitectónico y el control de
versiones, y se materializa mediante el desarrollo, las pruebas, el aseguramiento
de calidad y el despliegue. El mantenimiento, la gestión de configuración y la
mejora continua completan este ciclo y permiten conservar el valor de una
solución durante su evolución.

Sobre esta base, Legasoft se plantea como un espacio de trabajo colaborativo en
el que las decisiones funcionales, arquitectónicas y de calidad se mantienen
conectadas. La organización propuesta asigna responsabilidades claras a tres
roles compuestos, adecuados a un equipo en etapa inicial, y reconoce que la
comunicación permanente entre sus integrantes es necesaria para cumplir los
compromisos asumidos con cada cliente.

Este documento delimita los fundamentos institucionales, las líneas de servicio,
los objetivos, los referentes normativos y la estructura del equipo. La
infraestructura de trabajo, el acta que autoriza el inicio de actividades y las
hojas de vida de los integrantes se encuentran constituidas y se presentan como
anexos del documento de entrega \texttt{LGS-ENT-001}, del cual esta constitución
forma parte. La documentación técnica del proyecto —requisitos, arquitectura,
modelos y plan de pruebas— se elaborará durante el \lgSemestre\ sobre esa
infraestructura y siguiendo los procesos aquí definidos.

% =============================================================================
\section{Identidad corporativa de Legasoft}

\subsection{Origen y significado del nombre}

El nombre \emph{Legasoft} está inspirado en Valeri Legásov, científico soviético
que desempeñó un papel fundamental en la respuesta y el análisis del desastre de
Chernóbil. Su figura representa el compromiso con la verdad, la responsabilidad
profesional, la seguridad y el aprendizaje frente a los errores, principios que
la consultora busca aplicar en cada proyecto tecnológico. El Organismo
Internacional de Energía Atómica registra su participación como jefe de la
delegación soviética que presentó el análisis del accidente en 1986.

Asimismo, \emph{Lega} evoca la idea de legar o dejar un legado positivo mediante
soluciones que aporten valor duradero a las organizaciones, mientras que
\emph{Soft} hace referencia al software como principal medio para transformar
procesos, resolver problemas y generar innovación. En conjunto, Legasoft expresa
el propósito de crear tecnología responsable, confiable y capaz de dejar una
huella significativa en cada organización con la que colabora.

De ahí que la consultora no conciba la tecnología como un fin aislado, sino como
un medio para mejorar procesos, apoyar decisiones y generar capacidades
sostenibles. Un sistema que nadie puede mantener después de la entrega no se
considera un resultado satisfactorio.

\subsection{Perfil institucional}

Legasoft es una consultora tecnológica de carácter académico, en proceso de
formalización, especializada en transformación digital e ingeniería de software.
Su actividad se orienta a empresas, emprendimientos, instituciones y
organizaciones que requieren optimizar sus operaciones mediante soluciones
adaptadas a sus necesidades reales y a su capacidad de adopción tecnológica.

La intervención comienza con el análisis del funcionamiento actual del cliente.
Este diagnóstico permite reconocer procesos manuales, ineficiencias, dificultades
de gestión, restricciones y oportunidades de mejora. A partir de la evidencia
obtenida, la consultora documenta requisitos, propone una arquitectura apropiada
y desarrolla soluciones que favorecen la digitalización, la automatización y el
control de las operaciones.

El criterio institucional de Legasoft es mantener trazabilidad entre la necesidad
identificada, la solución diseñada y la calidad comprobada. Por ello, el trabajo
se apoya en documentación, control de versiones, pruebas, revisión técnica y
comunicación continua, con atención especial a la mantenibilidad, escalabilidad y
seguridad requeridas por el contexto de cada proyecto.

\subsection{Principio de intervención}

\begin{nota}[Principio rector]
Cada proyecto debe justificar la tecnología que incorpora.
\end{nota}

Legasoft prioriza soluciones proporcionales al problema y evita introducir
complejidad sin un beneficio verificable. Este principio guía la selección de
plataformas, patrones arquitectónicos, servicios en la nube, mecanismos de
automatización y componentes de inteligencia artificial o analítica de datos.

% =============================================================================
\section{Contexto organizacional y líneas de servicio}

Legasoft puede colaborar con pequeñas y medianas empresas, emprendimientos,
instituciones educativas, organizaciones sociales y áreas administrativas que
necesiten ordenar o modernizar su gestión. También considera organizaciones que
todavía dependen de tareas manuales, registros dispersos o sistemas existentes
que ya no responden adecuadamente a sus procesos.

La amplitud de este contexto no implica asumir dominio absoluto de todas las
tecnologías. La consultora define el alcance de cada intervención según el
diagnóstico, las competencias disponibles y los riesgos identificados. Cuando una
necesidad exceda su capacidad, deberá transparentarlo, ajustar el alcance o
gestionar el apoyo especializado correspondiente.

\subsection{Consultoría y transformación digital}

Esta línea se concentra en comprender cómo opera una organización antes de
proponer cambios tecnológicos. Incluye el diagnóstico de procesos, la
identificación de ineficiencias, la detección de tareas repetitivas y la
evaluación de oportunidades de digitalización o automatización. El resultado
esperado es una propuesta de mejora sustentada en necesidades comprobadas,
prioridades y restricciones del cliente.

La consultoría puede producir mapas de procesos, hallazgos, requisitos iniciales,
alternativas de solución y recomendaciones de implementación gradual. De este
modo, la inversión tecnológica se vincula con objetivos organizacionales
concretos y no únicamente con la incorporación de nuevas herramientas.

\subsection{Ingeniería y desarrollo de software}

Legasoft proyecta el diseño y desarrollo de sistemas web, aplicaciones móviles,
software a medida, sistemas administrativos, plataformas internas y sistemas de
gestión. El servicio comprende la definición del alcance, el análisis de
requisitos, el diseño de la solución, la implementación, la validación y la
preparación de la entrega.

El desarrollo se aborda como un proceso interdisciplinario. Las decisiones de
interfaz, lógica de negocio, persistencia de datos y despliegue deben responder a
los mismos requisitos y criterios de aceptación. Esta coordinación permite
construir productos con una base técnica comprensible y preparada para
mantenimiento posterior.

\subsection{Backend, APIs y datos}

Esta línea comprende el diseño e implementación de APIs, la construcción de
lógica de negocio, la integración con plataformas existentes y el modelado de
bases de datos. Cuando el contexto lo requiera, podrán emplearse gestores
relacionales como SQL Server, siempre que la selección se encuentre justificada
por las necesidades, restricciones y recursos del proyecto.

La gestión de información debe contemplar integridad, disponibilidad, acceso
autorizado y tratamiento responsable de los datos. Los contratos entre sistemas,
las reglas de negocio y las decisiones de modelado se documentarán para reducir
ambigüedades y facilitar la evolución de la solución.

\subsection{Experiencia de usuario y frontend}

Legasoft contempla el diseño de interfaces y el desarrollo frontend para entornos
web o móviles. El trabajo incluye organización de contenidos, construcción de
componentes, validación de flujos de interacción y adaptación de interfaces a
diferentes dispositivos. También puede abarcar la mejora de plataformas
existentes mediante componentes o widgets integrables.

La usabilidad y la accesibilidad se consideran atributos de calidad, no elementos
decorativos. Las decisiones visuales y de interacción deberán ayudar al usuario a
completar sus tareas con claridad, consistencia y una carga de aprendizaje
razonable.

\subsection{Cloud y DevOps}

Las prácticas de Cloud y DevOps apoyan el despliegue, la operación y la evolución
controlada del software. La consultora contempla contenerización, integración y
entrega continua, gestión de infraestructura en la nube, control de versiones y
mecanismos de disponibilidad, monitoreo y mantenimiento acordes con el alcance de
cada solución.

La adopción de estas prácticas será gradual y deberá considerar seguridad,
costos, capacidad operativa y continuidad. Su propósito es reducir errores
manuales, mejorar la repetibilidad de las entregas y conservar evidencia sobre
los cambios realizados.

\subsection{Inteligencia artificial y analítica}

Legasoft podrá incorporar inteligencia artificial, analítica de datos y
automatización inteligente únicamente cuando estas tecnologías aporten valor real
y exista información suficiente para utilizarlas de forma responsable. No se
presentarán como una respuesta predeterminada para cualquier problema.

Entre sus aplicaciones posibles se encuentran asistentes basados en inteligencia
artificial, recuperación y consulta de información, automatización de tareas con
criterios definidos, analítica descriptiva y apoyo a la toma de decisiones. Cada
incorporación deberá evaluar la calidad de los datos, la supervisión humana, la
seguridad, la privacidad y las limitaciones de los resultados obtenidos.

% =============================================================================
\section{Misión}

\begin{nota}[Misión]
Comprender las necesidades reales de cada organización y transformarlas en
soluciones de software personalizadas que permitan digitalizar y automatizar
procesos, mediante prácticas disciplinadas de ingeniería, comunicación continua y
aseguramiento de calidad, con el propósito de generar valor sostenible y
facilitar la evolución tecnológica del cliente.
\end{nota}

\section{Visión}

\begin{nota}[Visión]
Consolidarse como una consultora tecnológica reconocida por la calidad y
confiabilidad de su trabajo, la innovación responsable, la capacidad técnica de
su equipo y el cumplimiento de sus compromisos, desarrollando soluciones
mantenibles y escalables que acompañen de manera efectiva los objetivos de las
organizaciones atendidas.
\end{nota}

% =============================================================================
\section{Valores corporativos}

Los valores de Legasoft orientan tanto la relación con el cliente como las
decisiones internas del equipo. Su aplicación debe evidenciarse en la forma de
planificar, documentar, desarrollar, revisar y entregar cada solución.

\subsection{Valor y aplicación en Legasoft}

\begin{center}
\begin{tabularx}{\linewidth}{@{}L{3.6cm} Y@{}}
\toprule
\sffamily\bfseries\color{lgPrimario}Valor &
\sffamily\bfseries\color{lgPrimario}Aplicación en Legasoft \\
\midrule
Calidad & Se verifica que los entregables cumplan los requisitos, criterios de aceptación y atributos de calidad definidos para el proyecto. \\
Responsabilidad & Cada integrante asume sus compromisos, comunica oportunamente los riesgos y protege los recursos e información confiados por el cliente. \\
Transparencia & El estado del trabajo, las restricciones, las decisiones y las dificultades se informan de manera comprensible, sin generar expectativas que no puedan sostenerse. \\
Comunicación & Se mantienen canales claros para validar necesidades, coordinar actividades, resolver ambigüedades y conservar acuerdos documentados. \\
Innovación & Se exploran alternativas útiles y viables, evaluando su pertinencia, riesgos y aporte antes de incorporarlas a una solución. \\
Mejora continua & Las revisiones, métricas, defectos y lecciones aprendidas se emplean para ajustar los procesos y prevenir la repetición de problemas. \\
Orientación al cliente & Las prioridades técnicas se conectan con los objetivos, usuarios, restricciones y capacidad operativa de la organización atendida. \\
Trabajo en equipo & Los roles comparten información, revisan sus entregables y coordinan decisiones para conservar la coherencia integral del producto. \\
\bottomrule
\end{tabularx}
\captionof{table}{Valores corporativos y su aplicación práctica.}
\end{center}

% =============================================================================
\section{Objetivos de la consultora}

\subsection{Objetivo general}

Brindar servicios de transformación digital e ingeniería de software que
conviertan necesidades organizacionales verificadas en soluciones tecnológicas de
calidad, mantenibles y alineadas con los objetivos del cliente, mediante un
proceso documentado, colaborativo y sujeto a mejora continua.

\subsection{Objetivos específicos}

\begin{itemize}
  \item Diagnosticar los procesos y el contexto operativo de cada organización
        para identificar problemas, restricciones y oportunidades de mejora
        tecnológica.
  \item Levantar, analizar, priorizar y documentar requisitos con trazabilidad
        suficiente para orientar el diseño, el desarrollo y la validación de la
        solución.
  \item Diseñar arquitecturas y componentes tecnológicos proporcionales al
        alcance, la capacidad operativa y las necesidades reales del cliente.
  \item Desarrollar software mantenible, escalable y seguro mediante prácticas de
        control de versiones, revisión técnica y documentación de decisiones
        relevantes.
  \item Aplicar estrategias de pruebas y aseguramiento de calidad que permitan
        comprobar requisitos, criterios de aceptación y atributos del producto.
  \item Incorporar prácticas de Cloud y DevOps cuando contribuyan a despliegues
        repetibles, operación controlada y mantenimiento eficiente.
  \item Cumplir cronogramas y compromisos mediante seguimiento continuo, gestión
        de riesgos y revisión periódica de los procesos internos.
\end{itemize}

% =============================================================================
\section{Estándares y acreditaciones proyectadas}

Legasoft no cuenta actualmente con certificaciones empresariales formalmente
obtenidas. En su etapa inicial, la consultora busca alinear sus procesos con
normas reconocidas y utilizarlas como referencias para definir controles,
productos de trabajo y criterios de revisión. La aplicación será progresiva y
proporcional al tamaño del equipo y a la naturaleza de cada proyecto.

No todas las normas consideradas corresponden a certificaciones personales o a
certificaciones que la empresa deba obtener. Algunas describen modelos, procesos
o características que pueden emplearse como guía técnica. Como objetivo de
madurez organizacional, Legasoft podrá evaluar en el futuro esquemas de
certificación aplicables, una vez que existan procesos implementados, evidencia
suficiente y recursos para sostenerlos.

\begin{center}
\begin{tabularx}{\linewidth}{@{}L{3.4cm} Y@{}}
\toprule
\sffamily\bfseries\color{lgPrimario}Norma &
\sffamily\bfseries\color{lgPrimario}Uso previsto como referencia \\
\midrule
ISO 9001 & Sistema de gestión de la calidad y referencia para documentar procesos, responsabilidades, revisión de resultados y mejora continua. \\
ISO/IEC 27001 & Seguridad de la información y orientación para identificar riesgos y establecer controles sobre acceso, confidencialidad, integridad y disponibilidad. \\
ISO/IEC 12207 & Ciclo de vida del software y marco para organizar procesos de adquisición, desarrollo, operación, mantenimiento y apoyo del software. \\
ISO/IEC 25010 & Calidad del producto de software y modelo de características para definir y evaluar, según el proyecto, atributos como seguridad, usabilidad, fiabilidad y mantenibilidad. \\
ISO/IEC/IEEE 29148 & Ingeniería de requisitos y guía para levantar, analizar, especificar, validar y mantener requisitos con claridad y trazabilidad. \\
ISO/IEC 20000-1 & Gestión de servicios de TI y referencia futura para estructurar la prestación, seguimiento y mejora de servicios tecnológicos cuando la operación alcance mayor madurez. \\
\bottomrule
\end{tabularx}
\captionof{table}{Normas adoptadas como referencia, no como conformidad declarada.}
\end{center}

\begin{nota}[Alcance de esta declaración]
La adopción indicada no equivale a una declaración de conformidad ni a una
certificación obtenida. Todo avance deberá respaldarse con procesos definidos,
evidencia verificable, revisiones internas y, cuando corresponda, una evaluación
independiente realizada bajo el esquema aplicable.
\end{nota}

% =============================================================================
\section{Organización de la empresa}

Legasoft estará conformada inicialmente por tres integrantes con roles
compuestos. Esta estructura distribuye la coordinación, la arquitectura, el
desarrollo y la calidad sin crear niveles jerárquicos innecesarios. El Director
de Proyecto y Analista de Sistemas coordina el trabajo general y la relación con
el cliente; las decisiones técnicas y de calidad se construyen, revisan y
ejecutan de manera colaborativa.

\subsection{Organigrama inicial}

\begin{figure}[H]
\centering
\begin{tikzpicture}[
  font=\sffamily,
  cargo/.style={
    draw=lgPrimario, line width=0.9pt, rounded corners=3pt,
    fill=lgClaro, text width=5.0cm, align=center,
    inner sep=7pt, minimum height=1.9cm},
  % OJO: `align` debe ir después de `text width`; si no, TikZ vuelve a justificado.
  jefe/.style={cargo, fill=lgPrimario, text=white, draw=lgPrimario,
    text width=6.4cm, align=center},
  rotulo/.style={font=\sffamily\scriptsize\itshape, text=lgGris},
  coord/.style={-{Latex[length=2.4mm]}, lgPrimario, line width=0.9pt},
  inter/.style={{Latex[length=2.4mm]}-{Latex[length=2.4mm]}, lgAcento,
    line width=0.9pt, dashed}]

  \node[jefe] (dir)
    {\bfseries \lgDirectorRol\par\vspace{3pt}
     \normalfont\footnotesize \lgDirector};

  \node[cargo, below left=2.4cm and 0.35cm of dir.south] (arq)
    {\bfseries\color{lgPrimario} \lgArquitectoRol\par\vspace{3pt}
     \normalfont\footnotesize\color{black} \lgArquitecto};

  \node[cargo, below right=2.4cm and 0.35cm of dir.south] (qa)
    {\bfseries\color{lgPrimario} \lgQARol\par\vspace{3pt}
     \normalfont\footnotesize\color{black} \lgQA};

  % Coordinación descendente
  \coordinate (bus) at ($(dir.south)+(0,-1.2cm)$);
  \draw[lgPrimario, line width=0.9pt] (dir.south) -- (bus);
  \draw[lgPrimario, line width=0.9pt] (arq.north |- bus) -- (qa.north |- bus);
  \draw[coord] (arq.north |- bus) -- (arq.north);
  \draw[coord] (qa.north  |- bus) -- (qa.north);

  \node[rotulo, above=1pt of bus, anchor=south, fill=white, inner sep=2pt]
    {coordinación y seguimiento};

  % Intercambio técnico continuo
  \coordinate (busb) at ($(arq.south)+(0,-0.9cm)$);
  \draw[inter] (arq.south) -- (arq.south |- busb)
               -- (qa.south |- busb) -- (qa.south);
  \node[rotulo, anchor=north, yshift=-3pt]
    at ($(arq.south |- busb)!0.5!(qa.south |- busb)$)
    {intercambio técnico continuo};

\end{tikzpicture}
\caption{Organigrama inicial de Legasoft. Las líneas continuas representan la
coordinación general del proyecto; la línea discontinua, el intercambio
bidireccional entre arquitectura, backend, frontend y calidad.}
\end{figure}

El organigrama representa responsabilidades predominantes, no áreas aisladas. Los
tres integrantes participan en la planificación, las revisiones y la gestión de
cambios. Las flechas de coordinación muestran el seguimiento general del
proyecto, mientras que la relación bidireccional entre arquitectura, backend,
frontend y calidad refleja el intercambio continuo necesario para integrar el
producto.

% =============================================================================
\section{Definición de roles}

% ---------------------------------------------------------------- rol 1
\subsection{Director de Proyecto y Analista de Sistemas}

\begin{nota}[Titular del cargo]
\lgDirector
\end{nota}

\subsubsection{Propósito del cargo}

Conducir el proyecto desde la comprensión de la necesidad hasta la validación de
los resultados, manteniendo una relación clara con el cliente y una coordinación
efectiva del equipo. El cargo convierte objetivos organizacionales en requisitos
gestionables, organiza el trabajo y verifica que las decisiones del proyecto
conserven su alineación con el valor esperado.

\subsubsection{Responsabilidades}
\begin{itemize}
  \item Coordinar el proyecto y mantener la comunicación principal con los clientes.
  \item Planificar reuniones, entrevistas y actividades de levantamiento de información.
  \item Analizar procesos, problemas, necesidades, restricciones y prioridades.
  \item Elaborar y mantener la especificación de requisitos de software.
  \item Preparar minutas y conservar la documentación de las reuniones.
  \item Gestionar historias de usuario y acordar criterios de aceptación con el cliente.
  \item Organizar y supervisar el cronograma, las tareas y el seguimiento mediante Jira.
  \item Coordinar al equipo y facilitar la resolución de dependencias entre roles.
  \item Gestionar riesgos, solicitudes de cambio y prioridades del proyecto.
  \item Verificar que el producto y sus entregables respondan a los objetivos del cliente.
\end{itemize}

\subsubsection{Competencias técnicas}

Análisis de sistemas; levantamiento, especificación y gestión de requisitos;
documentación técnica y funcional; gestión ágil de proyectos; uso de Jira;
formulación de historias de usuario y criterios de aceptación; modelado de
procesos; aplicación de ISO/IEC/IEEE 29148 como referencia; conocimientos de
desarrollo web full stack; y conocimientos de Cloud y despliegue de aplicaciones.

\subsubsection{Competencias interpersonales}

Comunicación clara con clientes y equipo, escucha activa, facilitación de
reuniones, negociación de prioridades, liderazgo colaborativo, organización,
pensamiento analítico y capacidad para comunicar oportunamente riesgos o
restricciones.

\subsubsection{Entregables bajo su responsabilidad}
\begin{itemize}
  \item Plan de trabajo, cronograma y registro de seguimiento del proyecto.
  \item Minutas, acuerdos, registro de interesados y evidencia de validación.
  \item Especificación de requisitos, historias de usuario y criterios de aceptación.
  \item Registro de riesgos, cambios, decisiones funcionales y prioridades.
  \item Informes de avance y validación del cumplimiento de objetivos del cliente.
\end{itemize}

\subsubsection{Relación con los demás}

Trabaja con el Arquitecto de Software y Desarrollador Backend para transformar los
requisitos en decisiones técnicas viables y revisar el impacto de los cambios.
Coordina con el Especialista en QA y Desarrollador Frontend la definición de
criterios verificables, la validación de la experiencia de usuario y la
aceptación de los resultados. No sustituye el criterio técnico de estos roles;
integra sus aportes en la planificación y mantiene la trazabilidad con las
necesidades del cliente.

% ---------------------------------------------------------------- rol 2
\subsection{Arquitecto de Software y Desarrollador Backend}

\begin{nota}[Titular del cargo]
\lgArquitecto
\end{nota}

\subsubsection{Propósito del cargo}

Definir y materializar la estructura técnica de las soluciones, convirtiendo los
requisitos en una arquitectura coherente, segura y preparada para mantenimiento.
El cargo concentra el diseño de componentes, datos, integraciones y servicios
backend, y documenta las decisiones que condicionan la evolución del producto.

\subsubsection{Responsabilidades}
\begin{itemize}
  \item Diseñar la arquitectura técnica de las soluciones y representarla mediante el modelo C4.
  \item Seleccionar tecnologías y patrones de diseño de acuerdo con el contexto del proyecto.
  \item Diseñar el modelo de datos y utilizar SQL Server cuando corresponda.
  \item Diseñar e implementar APIs, integraciones y lógica de negocio.
  \item Definir y documentar contratos entre frontend y backend.
  \item Aplicar principios de seguridad en componentes, interfaces y acceso a datos.
  \item Documentar decisiones arquitectónicas y sus principales consecuencias.
  \item Revisar la escalabilidad, mantenibilidad e integridad técnica de la solución.
  \item Participar en integraciones, revisiones de código y actividades de despliegue.
  \item Colaborar en la estimación y análisis técnico de cambios o riesgos.
\end{itemize}

\subsubsection{Competencias técnicas}

Arquitectura de software; modelo C4; desarrollo backend; diseño de APIs REST;
bases de datos relacionales y SQL Server; modelado de datos; seguridad de
aplicaciones; patrones de diseño; control de versiones; principios SOLID;
integraciones; y documentación técnica de arquitectura.

\subsubsection{Competencias interpersonales}

Pensamiento sistémico, comunicación de decisiones técnicas, análisis de
alternativas, colaboración, disposición para la revisión entre pares,
responsabilidad sobre impactos técnicos y capacidad para explicar restricciones
sin perder de vista los objetivos del proyecto.

\subsubsection{Entregables bajo su responsabilidad}
\begin{itemize}
  \item Modelos de arquitectura C4 y documentación de decisiones arquitectónicas.
  \item Modelo de datos, definiciones de persistencia y scripts que correspondan al alcance.
  \item APIs, contratos de integración y componentes de lógica de negocio.
  \item Código backend revisado, versionado y acompañado por documentación técnica pertinente.
  \item Evidencia de revisiones de seguridad, mantenibilidad e integración técnica.
\end{itemize}

\subsubsection{Relación con los demás}

Valida con el Director de Proyecto y Analista de Sistemas que la arquitectura
responda a los requisitos, restricciones y prioridades acordados. Define con el
Especialista en QA y Desarrollador Frontend los contratos de integración, los
comportamientos verificables y las condiciones técnicas necesarias para las
pruebas. Participa en revisiones conjuntas para prevenir divergencias entre el
diseño, la implementación y la experiencia de usuario.

% ---------------------------------------------------------------- rol 3
\subsection{Especialista en QA y Desarrollador Frontend}

\begin{nota}[Titular del cargo]
\lgQA
\end{nota}

\subsubsection{Propósito del cargo}

Construir la interacción visible del producto y establecer controles de calidad
que permitan comprobar su comportamiento antes de cada entrega. El cargo integra
la perspectiva del usuario con la verificación sistemática de requisitos,
criterios de aceptación y atributos de calidad relevantes.

\subsubsection{Responsabilidades}
\begin{itemize}
  \item Desarrollar la interfaz de usuario para soluciones web o Flutter, según el proyecto.
  \item Implementar componentes y flujos coherentes con los requisitos y diseños acordados.
  \item Verificar la usabilidad y accesibilidad de las interfaces.
  \item Elaborar la estrategia de aseguramiento de calidad.
  \item Diseñar planes, escenarios y casos de prueba.
  \item Ejecutar pruebas funcionales y de integración dentro del alcance asignado.
  \item Registrar defectos, comunicar su impacto y realizar el seguimiento hasta su resolución.
  \item Verificar criterios de aceptación y validar la calidad antes de una entrega.
  \item Automatizar pruebas cuando resulte adecuado y sostenible para el proyecto.
  \item Participar en revisiones de código y apoyar procesos de integración y entrega continua.
\end{itemize}

\subsubsection{Competencias técnicas}

Desarrollo frontend y web; Flutter; diseño de interfaces; experiencia de usuario;
accesibilidad; pruebas funcionales y de integración; diseño de casos de prueba;
gestión de defectos; automatización de pruebas; CI/CD; control de versiones; y
aplicación de ISO/IEC 25010 como referencia para evaluar la calidad del producto.

\subsubsection{Competencias interpersonales}

Atención al detalle, pensamiento crítico, comunicación objetiva de defectos,
empatía con el usuario, organización de evidencia, colaboración con desarrollo,
criterio para priorizar riesgos y disposición para participar en ciclos de
retroalimentación.

\subsubsection{Entregables bajo su responsabilidad}
\begin{itemize}
  \item Interfaces y componentes frontend versionados y preparados para integración.
  \item Estrategia de calidad, plan de pruebas, escenarios y casos de prueba.
  \item Registro de ejecución, evidencias de prueba e informes de resultados.
  \item Registro y seguimiento de defectos con información suficiente para su reproducción.
  \item Informe de validación previo a la entrega y, cuando corresponda, pruebas automatizadas.
\end{itemize}

\subsubsection{Relación con los demás}

Coordina con el Director de Proyecto y Analista de Sistemas la interpretación de
los criterios de aceptación y comunica resultados que puedan afectar el alcance o
la fecha de entrega. Trabaja con el Arquitecto de Software y Desarrollador
Backend para acordar contratos, datos de prueba, integraciones y correcciones. Su
función de calidad se ejerce de manera preventiva desde la definición de
requisitos y no solamente al final del desarrollo.

% =============================================================================
\section{Presentación del equipo}

La composición inicial reúne capacidades de gestión, análisis, arquitectura,
desarrollo y calidad. Los roles principales delimitan la responsabilidad sobre
cada entregable, mientras que la integración del producto y las decisiones de
alcance permanecen como responsabilidades compartidas.

\begin{center}
\begin{tabularx}{\linewidth}{@{}L{4.0cm} L{3.4cm} Y@{}}
\toprule
\sffamily\bfseries\color{lgPrimario}Integrante &
\sffamily\bfseries\color{lgPrimario}Ámbito &
\sffamily\bfseries\color{lgPrimario}Aporte principal \\
\midrule
\lgDirector \newline \textit{\small \lgDirectorRol}
  & Gestión y requisitos
  & Coordinar el proyecto y asegurar la correspondencia entre necesidades, requisitos y resultados. \\
\lgArquitecto \newline \textit{\small \lgArquitectoRol}
  & Arquitectura y backend
  & Definir la estructura técnica e implementar datos, integraciones y lógica de negocio. \\
\lgQA \newline \textit{\small \lgQARol}
  & Frontend y calidad
  & Construir la interfaz y verificar el cumplimiento de los criterios de calidad y aceptación. \\
\bottomrule
\end{tabularx}
\captionof{table}{Composición inicial del equipo de Legasoft.}
\end{center}

% =============================================================================
\section{Hojas de vida del equipo}

Las hojas de vida de los tres integrantes, orientadas al rol que cada uno asume y
acompañadas de su disponibilidad horaria y su coste imputable al semestre, se
transcriben en el \textbf{Anexo C} del documento de entrega \texttt{LGS-ENT-001}.
Los archivos originales permanecen disponibles en los siguientes enlaces:

\begin{center}
\begin{tabularx}{\linewidth}{@{}L{4.6cm} Y@{}}
\toprule
\sffamily\bfseries\color{lgPrimario}Integrante y rol &
\sffamily\bfseries\color{lgPrimario}Hoja de vida \\
\midrule
\lgDirector \newline \textit{\small \lgDirectorRol}
  & \href{\lgDirectorCV}{\small Consultar documento} \\
\lgArquitecto \newline \textit{\small \lgArquitectoRol}
  & \href{\lgArquitectoCV}{\small Consultar documento} \\
\lgQA \newline \textit{\small \lgQARol}
  & \href{\lgQACV}{\small Consultar documento} \\
\bottomrule
\end{tabularx}
\captionof{table}{Enlaces a las hojas de vida del equipo.}
\end{center}

% =============================================================================
\section{Presencia y colaboración digital}

Legasoft dispone de un repositorio institucional en GitHub que funciona como
punto único de colaboración y de gestión de los activos de software y de
documentación de la consultora:

\begin{nota}[Repositorio institucional]
\href{\lgRepo}{\lgRepoCorto}
\end{nota}

El repositorio se encuentra creado y estructurado. Aloja la documentación
institucional y de proyecto, los modelos, el código fuente y las pruebas, y es
el medio por el cual se entrega y se verifica el avance del trabajo. Su
organización, las convenciones de control de versiones y el criterio con que se
distribuyen los entregables entre carpetas se detallan en el \textbf{Anexo A}
del documento de entrega \texttt{LGS-ENT-001}.

% =============================================================================
\section{Conclusiones}

La constitución inicial de Legasoft demuestra que una consultora de ingeniería de
software requiere responsabilidades definidas y coordinación efectiva, incluso
cuando está conformada por un equipo reducido. La distribución propuesta permite
identificar quién conduce la relación con el cliente y los requisitos, quién
resguarda la coherencia arquitectónica y del backend, y quién integra el frontend
con el aseguramiento de calidad. Esta claridad reduce vacíos de responsabilidad y
favorece decisiones compartidas con trazabilidad.

El valor del trabajo no depende solamente del código producido, sino de la
continuidad entre requisitos, arquitectura, desarrollo y calidad. Los procesos
documentados, el control de cambios, las pruebas y la comunicación permiten
comprobar que la solución responde al problema planteado y que puede mantenerse.
La adopción progresiva de normas como referencia ofrece una ruta de madurez, pero
no sustituye la necesidad de implementar prácticas y conservar evidencia antes de
declarar conformidad.

Legasoft busca ofrecer soluciones alineadas con necesidades reales, con una
innovación proporcional y responsable. Con esta versión la consultora completa su
base operativa: el acta que autoriza el inicio de actividades, la infraestructura
de trabajo con su repositorio estructurado y las hojas de vida del equipo quedan
formalizadas en los anexos de la entrega. Lo que sigue no es constituir la
empresa, sino ejercerla: seleccionar la organización cliente, levantar sus
requisitos y sostener a lo largo del \lgSemestre\ los procesos y la trazabilidad
que este documento se compromete a mantener.
